\scanpage{50}
\uline{Corollary} Every \deleted{\unsure{conti|cont}} saturated segment contains an onglet -- so if $\mathcal{R}$ is not the indiscrete relation there exist at least two onglets (for there then exist at least two saturated segments, complementary to each other)

Let \inserted{$\unsure{\Sigma_{1}|\Sigma_{0}}$ be the set of onglets of $R$,} $\deleted{\unsure{S_{0}|S_1}}\unsure{T_{1}|T_{0}}$ the union of the onglets, $\deleted{\unsure{S_1}}R_1$ its complement. Then, if $\deleted{S_1}R_1\ne\phi$, $\deleted{S_1}R_1$ has at least two component segments ($\operatorname{card}\unsure{\pi_0|\Pi_0}(R_1)\geq 2$), otherwise $R_1$ would be a \deleted{\unsure{s}}segment and would therefore contain an onglet, absurd. If $\operatorname{card} R_1=2$, then $R_1\in\mathcal{R}$ \deleted{is a} so $\mathcal{R}$ consists of the onglets and of $R_1$, (for \deleted{\unsure{$R_1$|$R$}} the $\unsure{S_i|S'}\subset R_1$ are of cardinality $\geq 2$, \deleted{\unsure{then|since}} for those of cardinality $1$ would be onglets \ldots)

If $\operatorname{card} R_1\geq 3$, $R_1$ is itself a polygonal graph; and $\mathcal{R}_1=\mathcal{R}\vert R_1$ is an unmixed equivalence relation. Moreover, \deleted{\unsure{al|Au}}No $\unsure{S_i|S_1}\in\mathcal{R}_1$ is contained in a component segment of $R_1$ -- otherwise it would be contained in a saturated segment contained in $\deleted{\unsure{S_1}}R_1$ (\unsure{Namely|namely} its cocomponent contained in this component), which would contain an onglet $\subset R_1$ -- absurd. If $\mathcal{R}_1$ is the indiscrete relation, we are done -- $\mathcal{R}$ is formed of the onglets, and of $R_1$.
% margin note written sideways
\margin{\unsure{Construction|Conclusion} \unsure{false}}
% archive sheet number: 49
